\chapter{Marked realization problems}
\label{v4-ch:marked-realizations}

An arithmetic construction can determine an algebra, a module, and a
trace without determining a geometric realization of them. A comparison
must specify which of these data it preserves. It must also specify the
maps used when the inputs vary. These requirements are already distinct
for two-dimensional modules over the dual numbers.

\section{A character does not determine a module}
\label{v4-marked:sec:characters}

Let \(k\) be a field of characteristic zero, and put
\(A=k[\epsilon]/(\epsilon^2)\). On the vector space \(k^2\),
with basis \(e_1,e_2\), consider the representations
\[
 \rho_0(a+b\epsilon)=aI,
 \qquad
 \rho_1(a+b\epsilon)=aI+bN,
 \qquad
 N=\begin{pmatrix}0&1\\0&0\end{pmatrix}.
\]
Both are unital algebra maps, since \(N^2=0\). Their characters,
the functions obtained by taking the matrix trace, agree on every
element of \(A\):
\begin{equation}\label{v4-marked:eq:characters}
 \operatorname{Tr}\rho_0(a+b\epsilon)
 =\operatorname{Tr}\rho_1(a+b\epsilon)=2a.
\end{equation}
They are nevertheless nonisomorphic as \(A\)-modules. An
intertwining isomorphism \(g\) would satisfy \(g\rho_0(\epsilon)
=\rho_1(\epsilon)g\), hence \(Ng=0\), contradicting
invertibility. This remains impossible if one permits an automorphism
of \(A\): such an automorphism sends \(\epsilon\) to
\(c\epsilon\) for some \(c\in k^\times\), so its nonzero
nilpotent action remains nonzero.

The vector \(e_2\) gives a further useful mark. It generates the
whole module under \(\rho_1\), since \(Ne_2=e_1\); under
\(\rho_0\) it generates only \(ke_2\). The common character has
forgotten an extension and a cyclic-vector condition. Even equality of
all these scalar evaluations does not provide the module comparison.

\section{The data carried by a comparison}
\label{v4-marked:sec:data}

Fix a commutative ring \(R\). Complexes are cohomological and all
tensor products in this section are over \(R\). A \emph{marked
algebra with module and functional} is a tuple
\[
 \mathcal A=(A,M,v,\tau),
\]
where \(A\) is a unital associative differential graded algebra,
\(M\) is a left differential graded \(A\)-module,
\(v:R\to M\) is a chain map, and \(\tau:A\to R\) is a
chain map of degree zero. The chosen vector \(v(1)\) is therefore
closed. If the functional is required to be cyclic, this means
\[
 \tau(ab)=(-1)^{|a||b|}\tau(ba)
\]
for homogeneous inputs. Cyclicity is an extra hypothesis; it is not
part of the meaning of an arbitrary functional.

A comparison \((f,g):\mathcal A\to\mathcal A'\) consists of a
unital differential graded algebra map \(f:A\to A'\) and a
chain map \(g:M\to M'\) such that
\begin{equation}\label{v4-marked:eq:comparison}
 g(am)=f(a)g(m),\qquad
 g\circ v=v',\qquad
 \tau'\circ f=\tau.
\end{equation}
An isomorphism requires both maps to be invertible. A comparison need
not be an isomorphism: its image and the operations lost in its kernel
can be part of the question.

The two distinguished maps in this tuple have different types. The
algebra unit is a map \(R\to A\). The selected vector is a map
\(R\to M\). Unitality of the action says \(1_Am=m\); it does
not choose \(v\). In geometric constructions an empty coefficient
and a local vacuum likewise need their own maps on their respective
carriers.

Additional operations enlarge the tuple and the equations
\eqref{v4-marked:eq:comparison}. For example, an endomorphism
\(F_A:A\to A\) is preserved when \(fF_A=F_{A'}f\); an
operator \(F_M\) on the module requires \(gF_M=F_{M'}g\).
An involution requires its own compatibility. If a multiplication is
associative only through specified higher homotopies, the comparison
contains components compatible with those homotopies. Keeping only
its first component is then a different comparison problem.

\begin{proposition}\label{v4-marked:prop:composition}
The displayed marked objects and their comparisons form a category.
Every comparison preserves an ordered module operation and its final
evaluation whenever the latter is defined by the retained maps. In
particular,
\[
 g(a_1\cdots a_r v(1))
 =f(a_1)\cdots f(a_r)v'(1),
 \qquad
 \tau'(f(a_1)\cdots f(a_r))=\tau(a_1\cdots a_r).
\]
\end{proposition}

\begin{proof}
Compose the algebra maps and the module maps separately. Substitution
in each equation of \eqref{v4-marked:eq:comparison} verifies it for
the composite; identity maps satisfy the same equations. Ordinary
composition of chain maps gives associativity. The first displayed
identity follows by induction using the module equation. The second
uses multiplicativity of \(f\) and then \(\tau'f=\tau\).
No factors have been interchanged, so no new Koszul sign occurs.
\end{proof}

The converse is false if only scalar identities are retained. Equation
\eqref{v4-marked:eq:characters} gives characters agreeing on every
input while the modules are not isomorphic. The ordered matrix calculation
in the prologue shows why even a scalar identity at one arity may fail
to survive insertion into a larger operation.

\section{Coefficient change and geometric carriers}
\label{v4-marked:sec:coefficients}

A change of coefficient ring begins with an actual unital map
\(R\to S\) of commutative rings. Ordinary extension gives
\[
 A_S=S\otimes_R A,\qquad M_S=S\otimes_R M,
 \qquad v_S=1\otimes v,\qquad \tau_S=1\otimes\tau,
\]
using \(S\otimes_R R\cong S\). Its algebra product and action
are the induced ones. Tensoring the equations
\eqref{v4-marked:eq:comparison} gives a comparison over \(S\).
This algebraic construction does not require flatness. Extension
depending only on quasi-isomorphism classes uses derived tensor
products, with compatible algebra and module resolutions when needed.

For a prime \(\ell\) and \(r\ge1\), there is no unital ring
map \(\mathbb Z/\ell^r\to K\) to a field \(K\) of
characteristic zero: the relation \(\ell^r1=0\) would have to
hold in \(K\). More generally, if an abelian group \(V\) is
annihilated by \(\ell^r\), then
\[
 V\otimes_{\mathbb Z}K=0,
\]
because
\(v\otimes c=\ell^rv\otimes(c/\ell^r)=0\).
Thus direct rationalization of finite-torsion cochains loses their
coefficients.

The inverse-limit construction is different. With the usual reduction
maps,
\[
 \mathbb Z_\ell=\varprojlim_r\mathbb Z/\ell^r,
 \qquad \mathbb Q_\ell=\mathbb Z_\ell[1/\ell].
\]
Every term of the inverse system tensors to zero over \(\mathbb Z\)
with \(\mathbb Q\), whereas
\(\mathbb Z_\ell\otimes_{\mathbb Z}\mathbb Q
\cong\mathbb Q_\ell\ne0\): an integer prime to \(\ell\)
is already a unit in \(\mathbb Z_\ell\). Tensoring with
\(\mathbb Q\) has therefore not commuted with this inverse
limit. There is also no continuous unital field map
\(\mathbb Q_\ell\to\mathbb C\) with their usual topologies.
Such a map fixes the rational numbers, but the sequence \(\ell^n\)
tends to zero in \(\mathbb Q_\ell\) and has unbounded complex
absolute value. This argument concerns continuity, not the existence
of an abstract field embedding.

Geometry introduces another change of carrier. If an arithmetic input
is defined over \(R\) and the position curve \(X\) is defined
over a field \(k\), the coefficient functor and its interaction with the
operations must be specified. Even after a scalar map \(R\to k\)
has been chosen, a coefficient algebra alone gives no sheaves on
the Cartesian powers \(X^I\), for finite ordered sets \(I\), no
restrictions to collision strata, and no maps for
inserting one configuration into another.

For a chosen family of sheaves \(A_I\), write
\(\mathcal Q_n(A_{I_1},\ldots,A_{I_n};A_J)\) for its declared
complex of allowed operations. The definition of \(\mathcal Q_n\)
must specify supports, pole or growth bounds, tensor products and
completions, as appropriate to the example. A geometric comparison
is then a system of maps on these actual complexes preserving their
differentials, substitutions and specified units. A primitive in a
larger Hom complex need not belong to \(\mathcal Q_n\).
Consequently an obstruction being exact after forgetting support or
continuity does not provide a comparison in the chosen category.

\section{Existence, uniqueness, and coherent choice}
\label{v4-marked:sec:spaces}

Suppose that a category of realizations and a category of retained
data have been constructed, together with a forgetful functor.
A groupoid is a category in which every arrow is invertible.
To ask for realizations up to isomorphism, retain all objects and
only invertible arrows. Denote the resulting groupoids by
\(\mathscr R\) and \(\mathscr D\), and the functor between
them by \(U:\mathscr R\to\mathscr D\).
For \(d\in\mathscr D\), its \emph{realization groupoid}
has objects
\begin{equation}\label{v4-marked:eq:fibre}
 \operatorname{Ob}\operatorname{Real}_U(d)
 =\{(r,\alpha):r\in\mathscr R,
                 \ \alpha:U(r)\xrightarrow{\sim}d\}.
\end{equation}
An arrow from \((r,\alpha)\) to \((r',\alpha')\) is an
arrow \(h:r\to r'\) satisfying \(\alpha' U(h)=\alpha\).
This description is the homotopy fibre of the map of classifying
spaces of the two groupoids.

Existence means that this groupoid has an object. Uniqueness up to
isomorphism means that it is connected. Uniqueness with a unique
compatible isomorphism means that it is equivalent to the one-object
groupoid with only its identity arrow; its classifying space is then
contractible. Automorphisms of a realization can prevent this last
condition even when every realization is isomorphic to it.

In a setting with higher morphisms the analogous definition is
\[
 \operatorname{Real}_U(d)
 =\operatorname{hofib}_d
   (U^\simeq:\mathscr R^\simeq\to\mathscr D^\simeq),
\]
where the superscript denotes the space retaining objects and all
equivalences. That formula presupposes the categories and \(U\);
it does not construct them. A comparison required to be noninvertible
belongs to the relevant category of arrows, with its prescribed
source and target, rather than to the displayed space of
identifications.

A functor is fully faithful when its maps on all sets of arrows are
bijections. It is essentially surjective when every target object
is isomorphic to an image object.

\begin{proposition}\label{v4-marked:prop:contractible-fibres}
For a functor between groupoids, every realization groupoid is
equivalent to a point if and only if the functor is fully faithful
and essentially surjective.
\end{proposition}

\begin{proof}
Essential surjectivity supplies an object \((r,\alpha)\) in
each fibre. Full faithfulness supplies the unique arrow between any
two such objects, namely the lift of \((\alpha')^{-1}\alpha\).
This makes each fibre equivalent to a point.

Conversely, nonemptiness of the fibres gives essential surjectivity.
Given \(\beta:U(r)\to U(r')\), consider the two objects
\((r,\beta)\) and \((r',1)\) in the fibre over \(U(r')\).
Its equivalence to a point supplies exactly one arrow between them.
By its defining equation, this is exactly one lift of \(\beta\).
Thus \(U\) is fully faithful.
\end{proof}

Objectwise existence alone does not give coherent choices. For an
explicit example, let \(C_n=\mathbb Z/n\mathbb Z\), written
additively, and let \(BC_n\) be the groupoid with one object and
automorphism group \(C_n\). Reduction modulo two gives
\[
 U:BC_4\longrightarrow BC_2.
\]
Its realization groupoid has two objects, corresponding to the two
identifications in \(C_2\). Reduction \(C_4\to C_2\) is
surjective, so they are isomorphic. The automorphism group of either
is the kernel \(\{0,2\}\cong C_2\). Thus the fibre is
equivalent to \(BC_2\); in particular it is nonempty.

A functorial section, even one whose composite with \(U\) is
only naturally isomorphic to the identity, would give a group
homomorphism \(s:C_2\to C_4\) splitting reduction. The natural
isomorphism changes the composite by conjugation, which is trivial
in \(C_2\). But the image of the nonzero element would have to
have order dividing two, so it is \(0\) or \(2\); both reduce
to zero. No such section exists. The missing compatibility is on
arrows, although the objectwise existence problem has been solved.

\section{Evaluation after the comparison}
\label{v4-marked:sec:evaluation}

The functional equation in \eqref{v4-marked:eq:comparison} is a
commutative triangle on a common domain. If the two functionals are
defined independently, its proof identifies their evaluations.
If instead \(\tau\) is defined to be \(\tau'f\), the same
triangle holds by definition and gives a transported functional.
It has not identified an independently prescribed arithmetic one.
Both constructions are legitimate, with different conclusions.

For positivity the algebraic structure being preserved is decisive.
Let \(A,B\) be complex algebras with conjugate-linear involutions
\(*\) satisfying \((ab)^*=b^*a^*\), and let
\(f:A\to B\) preserve the product and involution. Suppose that
\(\tau_B\) is linear on a subspace containing every
\(f(a)^*f(a)\), and
\(\tau_B(f(a)^*f(a))\ge0\) for all \(a\in A\).
If \(\tau_A(a^*a)=\tau_B(f(a^*a))\) on the same inputs, then
\[
 \tau_A(a^*a)=\tau_B(f(a)^*f(a))\ge0.
\]
This implication uses the actual product and involution identities.
For an unbounded weight, the displayed expressions must lie in its
domain; no value at the unit is inferred. A proposed positive trace
realization of the Weil functional must supply this map and this
evaluation identity on its full test domain.

The construction of each side is part of the comparison problem.
Two arguments may use the same homological algebra without being
circular; disjoint lists of references can still conceal a target
defined from the desired answer. Independence here requires the
arithmetic input and the comparison object to be specified without
defining one by the invariant that the comparison is meant to
explain. The remaining task is to construct the operation-preserving
maps and their coherent behavior under change of inputs.
